\scanpage{14}
\mspage{5}
% archive sheet number: 14

\deleted{Tout ceci est récapitulatif, ne fait pas \deleted{\unsure{intervenir}}\inserted{intervenir} directement $\mathrel{|\unsure{\circ|0}|}$.}

\deleted{\uline{Prop. 4}}\deleted{\uline{Prop \unsure{préc}}}
\margin{devrait \unsure{remonter} après prop 3.}
% slanted

\uline{Cor 2} Soit $G\ll F=\bigcup F_i$. Alors $G\leq F$ ssi $\forall i,\ G|F_i\leq F_i$. \deleted{\unsure{Alors}} Si $G,G'\ll F$, alors $G\unsure{\gtrless|\bowtie}G'$ ssi \inserted{$\forall i$} $G_{F_i}\gtrless G'_{F_i}$, idem pour $\parallel$, $\leq$, $\ll$ \ldots

Tout ceci est récapitulatif, et ne fait pas intervenir explicitement $\mathrel{|\circ|}$, si ce n'est par l'intermédiaire de $\gtrless$. Il faudrait, dans ce contexte, revoir \unsure{aussi|assez} rapidement

1) Les axiomes relatifs aux Lieux, ombres, multiombres.

2) Le \unsure{lien} avec le cas ensembliste et quasi-ensembliste (où les lieux jouent \unsure{justement} un rôle prépondérant).

Mais surtout

3) Revoir la théorie des supports.

\deleted{\unsure{Comme}} Également

4) Voir ce que devient l'axiomatique $\leq$, $\ll$, $\mathrel{|\circ|}$, dans le cas (le plus courant sûrement) où $\ll$ peut s'exprimer en termes de $\leq$, $\mathrel{|\circ|}$. Mais pour le formuler, nous \unsure{aurons|avons} besoin de donner un sens, pour $X\in M$, à
\[
\text{``}\ X^\circ=X\setminus\partial X\ \text{''},
\]
et pour $X,Y\in M$, à ``l'inclusion''
\[
\text{``}\ X^\circ\subset Y^{\unsure{\bullet|\circ}}\ \text{''}\qquad\text{---}
\]
i.e. il faut d'abord passer par 3.

Traitons quand même $1^\circ)$ et $2^\circ)$.

\deleted{Les axiomes \unsure{pertinents deviennent}}
